\section{Introducción: RAG es un problema de sistemas, no «buscar y luego armar el Prompt»}

La conferencia de Douwe Kiela parte de la historia de los modelos de lenguaje y va agregando, capa por capa, retrieval, reranking, fusion, joint training, active retrieval y tool use. El eje de toda la clase no es un pipeline convertido en producto, sino una pregunta de investigación: cómo entra el conocimiento externo al proceso de generación, qué componentes deben congelarse o entrenarse, cuándo se recupera la evidencia, cómo se ordena, dónde se fusiona y cómo demuestra el sistema que su salida está realmente restringida por la evidencia.

\lecturefigure{slide-01-title.jpg}{Página de título de la clase Retrieval Augmented Language Models}{Video de la clase de Stanford 00:01:39}

\begin{knowledgebox}{Lectura de la figura: el Augmented del título cambia los límites del modelo}
Un language model común concentra la mayor parte de su conocimiento en sus parámetros; un retrieval-augmented language model accede, durante la inferencia o el entrenamiento, a documentos externos, índices vectoriales, motores de búsqueda u otras herramientas. Lo que se amplía no es solo la longitud de la entrada, sino la interfaz entre el modelo y una non-parametric memory que puede actualizarse.
\end{knowledgebox}

\begin{importantbox}{La pregunta común de esta clase}
Sea $x$ la entrada, $\mathcal{D}$ el corpus externo, $R$ el recuperador y $G$ el generador; el flujo más simple es
\[
z_{1:k}=R(x,\mathcal{D}),\qquad
y=G(x,z_{1:k}).
\]
La verdadera dificultad ocurre dentro de las dos igualdades: cómo aprende $R$ a obtener evidencia «útil para $G$», si $G$ realmente usará $z_{1:k}$ y cómo equilibra el sistema completo calidad, costo, actualización y confiabilidad.
\end{importantbox}

\subsection{Límites de la evidencia y método de lectura}

Esta clase se impartió en diciembre de 2023. Las valoraciones del ponente sobre DRAGON, vector database, retrieval hardware y la estructura futura de la industria se conservan como opiniones situadas en la fecha de la clase; los mecanismos de RAG, REALM, FiD, RETRO, Atlas, RePlug y otros se explican con base en los artículos originales. El dashboard, el incident runbook, el gate dropout, el drift monitor y la gobernanza regional que aparecían en la versión anterior no tienen respaldo en la clase, por lo que se retiraron por completo.

\subsection{Resumen de la sección}

La línea didáctica de RAG va de «documentos externos que entran al prompt» a «recuperación, representación, fusión y generación optimizadas en conjunto». Las secciones siguientes explican primero por qué se necesita una memoria externa, luego comparan arquitecturas con la train/test taxonomy y, por último, analizan la dirección de sistematización de RAG 2.0.
